\chapter{Đầu ra nhanh duy nhất}
% full version: chapters/13_muscles.tex

Mọi việc bạn làm nhanh trong thế giới --- một lời nói, một ánh nhìn, một bước chân --- đều là co cơ. Cỗ máy còn có một đầu ra thứ hai, chậm, bằng hóa chất, nói ở chương mười sáu. Còn ở đây --- cơ bắp.

Mắt xích cuối cùng của chuỗi là các nơ-ron \nt{ACET}. Mỗi sợi cơ nghe đúng một nơ-ron như vậy, còn một nơ-ron điều khiển một nhóm sợi, từ vài trăm đến vài nghìn. Điều thú vị là mối nối giữa dây thần kinh và cơ là mối nối đáng tin cậy nhất trong cơ thể. Bên trong não, các mối nối làm mất hơn một nửa số tín hiệu, còn ở đây mỗi tiếng tách nhả ra lượng chất nhiều gấp mấy lần mức cần thiết. Ở đầu ra, một lỗi phải trả giá bằng cả một cử động, nên độ tin cậy được mua dư dả.

\section*{Hồi trống thành cú đẩy}

Một tiếng tách --- một lần giật cơ ngắn. Nếu tiếng tách đến thưa, cơ giật từng cái riêng lẻ. Nếu đến dày, các lần giật hòa vào nhau thành một lực đều, như hồi trống dồn hòa thành tiếng ù. Cơ biến dòng xung thành lực, như một bộ chuyển đổi số--tương tự đơn giản nhất.

\pencilfig{113}{%
% twitch = alpha function; the sum of twitches for spikes 1 s and 0.16 s apart (arbitrary units of time)
\begin{scope}[x=0.95cm, y=0.36cm]
\foreach \dx/\t in {0/thưa,4.6/dày}{
  \begin{scope}[xshift=\dx cm]
  \draw[pen] (0,0) -- (4,0); \node[font=\small\itshape, text=pcGray, anchor=south] at (2,5.4) {\t};
  \end{scope}}
\draw[penlite=pcRed] plot[smooth] coordinates {(0.00,0.00) (0.05,0.00) (0.10,0.00) (0.15,0.00) (0.20,0.00) (0.25,0.45) (0.30,0.73) (0.35,0.90) (0.40,0.98) (0.45,1.00) (0.50,0.98) (0.55,0.94) (0.60,0.88) (0.65,0.81) (0.70,0.74) (0.75,0.66) (0.80,0.59) (0.85,0.52) (0.90,0.46) (0.95,0.41) (1.00,0.35) (1.05,0.31) (1.10,0.27) (1.15,0.23) (1.20,0.20) (1.25,0.62) (1.30,0.88) (1.35,1.02) (1.40,1.08) (1.45,1.09) (1.50,1.06) (1.55,1.00) (1.60,0.93) (1.65,0.86) (1.70,0.78) (1.75,0.70) (1.80,0.62) (1.85,0.55) (1.90,0.48) (1.95,0.42) (2.00,0.37) (2.05,0.32) (2.10,0.28) (2.15,0.24) (2.20,0.21) (2.25,0.62) (2.30,0.88) (2.35,1.03) (2.40,1.09) (2.45,1.09) (2.50,1.06) (2.55,1.01) (2.60,0.94) (2.65,0.86) (2.70,0.78) (2.75,0.70) (2.80,0.62) (2.85,0.55) (2.90,0.48) (2.95,0.42) (3.00,0.37) (3.05,0.32) (3.10,0.28) (3.15,0.24) (3.20,0.21) (3.25,0.62) (3.30,0.88) (3.35,1.03) (3.40,1.09) (3.45,1.09) (3.50,1.06) (3.55,1.01) (3.60,0.94) (3.65,0.86) (3.70,0.78) (3.75,0.70) (3.80,0.62) (3.85,0.55) (3.90,0.48) (3.95,0.42) (4.00,0.37)};
\foreach \s in {0.20,1.20,2.20,3.20}{\draw[pcBlue, line width=0.8pt] (\s,-0.35) -- (\s,-1.2);}
\begin{scope}[xshift=4.6cm]
\draw[penlite=pcRed] plot[smooth] coordinates {(0.00,0.00) (0.05,0.00) (0.10,0.00) (0.15,0.00) (0.20,0.00) (0.25,0.45) (0.30,0.73) (0.35,0.90) (0.40,1.35) (0.45,1.68) (0.50,1.85) (0.55,2.19) (0.60,2.51) (0.65,2.64) (0.70,2.84) (0.75,3.13) (0.80,3.22) (0.85,3.27) (0.90,3.55) (0.95,3.62) (1.00,3.54) (1.05,3.82) (1.10,3.88) (1.15,3.80) (1.20,3.98) (1.25,4.06) (1.30,3.97) (1.35,4.07) (1.40,4.16) (1.45,4.09) (1.50,4.10) (1.55,4.23) (1.60,4.17) (1.65,4.10) (1.70,4.26) (1.75,4.23) (1.80,4.07) (1.85,4.27) (1.90,4.27) (1.95,4.12) (2.00,4.26) (2.05,4.29) (2.10,4.17) (2.15,4.24) (2.20,4.31) (2.25,4.21) (2.30,4.21) (2.35,4.31) (2.40,4.24) (2.45,4.16) (2.50,4.31) (2.55,4.27) (2.60,4.10) (2.65,4.30) (2.70,4.29) (2.75,4.15) (2.80,4.28) (2.85,4.31) (2.90,4.18) (2.95,4.25) (3.00,4.32) (3.05,4.22) (3.10,4.21) (3.15,4.32) (3.20,4.25) (3.25,4.17) (3.30,4.31) (3.35,4.27) (3.40,4.11) (3.45,3.86) (3.50,3.56) (3.55,3.25) (3.60,2.93) (3.65,2.63) (3.70,2.33) (3.75,2.06) (3.80,1.81) (3.85,1.58) (3.90,1.38) (3.95,1.19) (4.00,1.03)};
\foreach \s in {0.20,0.36,0.52,0.68,0.84,1.00,1.16,1.32,1.48,1.64,1.80,1.96,2.12,2.28,2.44,2.60,2.76,2.92,3.08,3.24}{\draw[pcBlue, line width=0.8pt] (\s,-0.35) -- (\s,-1.2);}
\end{scope}
\node[lbl, text=pcRed, anchor=west] at (1.2,1.9) {các lần giật};
\node[lbl, text=pcRed, anchor=south] at (6.6,4.55) {lực đều};
\node[lbl, text=pcBlue, anchor=north] at (4.3,-1.35) {tiếng tách của nơ-ron};
\end{scope}
}{Tiếng tách thưa cho những lần giật cơ riêng lẻ; tiếng tách dày hòa thành một lực đều và lớn hơn nhiều.}

\section*{Lực dấu phẩy động}

Bộ não có hai núm chỉnh lực: nơ-ron tách dày đến đâu và bao nhiêu nhóm sợi cơ được bật. Các nhóm bật lên theo thứ tự --- từ yếu nhất tới mạnh nhất, và không ai tính toán thứ tự này: các nơ-ron nhỏ đơn giản là dễ bị kích thích hơn. Kết quả là mỗi nhóm mới thêm vào gần như cùng một \emph{tỷ lệ} so với lực đã có. Khi bạn cầm quả trứng, lực được định lượng bằng những phần cực nhỏ, khi bạn nhấc va li --- bằng những phần lớn. Kỹ sư sẽ nhận ra ở đây một số dấu phẩy động: bậc độ lớn do số nhóm đã bật quyết định, còn giá trị chính xác do tần số tiếng tách. Mặt trái: cử động càng mạnh, càng kém chính xác. Theo một giả thuyết có ảnh hưởng, đó là lý do các cử động của chúng ta mượt như vậy: lệnh mượt cho độ tản mát nhỏ nhất.

\section*{Kéo chứ không đẩy}

Cơ chỉ biết kéo. Vì thế mỗi khớp được phục vụ bởi một cặp cơ kéo về hai phía ngược nhau --- lại là hai dây dẫn thay cho dấu. Hiệu số lực của chúng xoay khớp, còn tổng làm khớp cứng hơn. Khi bạn bưng một tách đầy, bạn gồng cả hai cơ cùng lúc: góc xoay vẫn thế, nhưng cánh tay chịu va chạm vững hơn.

\pencilfig{13}{\pencilart{13}}{Cơ chỉ biết kéo, nên có hai cơ: hiệu số lực xoay khớp, tổng làm khớp vững hơn.}

\section*{Phản hồi bị trễ}

Trong cơ có các cảm biến căng, và vòng lặp ngắn nhất khép kín ngay trong tủy sống: cơ bị kéo căng --- nó tự co lại, còn cơ đối kháng tạm thời bị tắt. Thứ tắt nó là một nơ-ron ức chế loại \nt{GLYC}: loại này đã có việc để làm.

Nhưng mỗi vòng lặp đều có độ trễ --- trong tủy sống là vài chục mili giây, qua mắt là vài trăm. Mà phản hồi bị trễ làm hệ thống dao động, như người lái xe cầm lái mà nhìn đường chậm một nhịp. Vòng lặp càng dài thì càng phải sửa chậm. Ngưỡng mà vòng lặp tủy sống bắt đầu dao động là khoảng tám dao động mỗi giây. Con số này gần với tần số của cơn run nhẹ mà người khỏe mạnh nào cũng có, và vòng lặp căng cơ là một trong những nguồn gốc khả dĩ của nó.

Vậy làm sao chúng ta cử động được nhanh và chính xác? Theo một giả thuyết phổ biến --- bằng dự đoán: giữ một mô hình bên trong của cơ thể và tính trước kết quả của lệnh, không chờ các cảm biến.

\section*{Nhịp không cần máy đếm nhịp}

Đi bộ và hít thở đều có nhịp, nhưng chúng không cần máy đếm nhịp. Hai nhóm nơ-ron ức chế lẫn nhau và mỏi dần sẽ tự thay phiên nhau làm việc: kẻ thắng mệt đi, đối thủ đã nghỉ ngơi giành phần hơn, và cứ thế mãi. Một lệnh không đổi từ trên xuống --- ``đi'' --- ngay tại chỗ biến thành nhịp ``trái-phải''.

\fullversion{13}
